% Presentación corta de VacuInfo
% Compilar con: latexmk -pdf vacuinfo.tex
\documentclass[aspectratio=169,11pt]{beamer}

\usepackage[utf8]{inputenc}
\usepackage[T1]{fontenc}
\usepackage[spanish,es-noshorthands]{babel}
\usepackage{booktabs}
\usepackage{tikz}
\usetikzlibrary{positioning,arrows.meta}

\usetheme{Madrid}
\setbeamertemplate{navigation symbols}{}
\setbeamerfont{footline}{size=\scriptsize}

\definecolor{azul}{HTML}{1F4E79}
\definecolor{azulclaro}{HTML}{DEEAF6}
\definecolor{naranja}{HTML}{C55A11}
\definecolor{naranjaclaro}{HTML}{FBE5D6}
\definecolor{grisclaro}{HTML}{F2F2F2}

\title[NTS 246]{Consulta de la normativa de vacunación del Perú con RAG}
\author[Grupo 6]{%
  \begin{tabular}{c}
    Juan Chipana \quad Ximena Quispe \quad Sean Leiva \\[3pt]
    Marcela Ventura \quad Dylan Lucar \quad Alejandro Ramirez
  \end{tabular}}
\date{Noviembre 2026}

\tikzset{
  caja/.style={
    draw=azul, fill=azulclaro, rounded corners=2pt,
    text width=1.85cm, minimum height=1.0cm, align=center,
    font=\scriptsize, inner sep=3pt
  },
  cajaN/.style={
    draw=naranja, fill=naranjaclaro, rounded corners=2pt,
    text width=1.85cm, minimum height=1.0cm, align=center,
    font=\scriptsize, inner sep=3pt
  },
  arch/.style={
    draw=azul, fill=grisclaro, rounded corners=2pt,
    text width=1.85cm, minimum height=1.0cm, align=center,
    font=\scriptsize, inner sep=3pt
  },
  flecha/.style={->, thick, azul},
  flechaN/.style={->, thick, naranja}
}

\begin{document}

% ---------------------------------------------------------------- Portada
\begin{frame}
  \titlepage
\end{frame}

% ---------------------------------------------------------------- Objetivo
\begin{frame}{Objetivo}
  \begin{block}{Objetivo general}
    Construir un sistema de \textbf{RAG (Retrieval-Augmented Generation)} que permita
    consultar en lenguaje natural la \textbf{NTS N° 246-MINSA/DGIESP-2026},
    \emph{Norma Técnica de Salud que establece el Esquema Nacional de Inmunizaciones}.
  \end{block}

  \vspace{2mm}
  \begin{itemize}
    \item Responder \textbf{únicamente} con base en los PDF oficiales de \texttt{pdf/}.
    \item Devolver la \textbf{página y sección} que sustentan cada respuesta.
    \item Ofrecer una consulta \textbf{rápida y trazable} frente a la búsqueda manual.
    \item Ser una herramienta de consulta, \textbf{no} un asistente médico ni un reemplazo
          de la lectura de la norma.
  \end{itemize}
\end{frame}

% ---------------------------------------------------------------- Problema
\begin{frame}{El problema}
  \begin{columns}[T]
    \begin{column}{0.52\textwidth}
      \begin{itemize}
        \item Norma de \textbf{140 páginas escaneadas} (97 + 43), con numeración
              impresa continua.
        \item Organizada en secciones, fichas técnicas de vacunas y tablas.
        \item Buscar una edad, dosis, intervalo, temperatura de conservación o
              contraindicación exige revisar muchas páginas.
        \item Los PDF \textbf{no tienen capa de texto}: hay que aplicar OCR antes
              de poder buscar.
      \end{itemize}
    \end{column}
    \begin{column}{0.44\textwidth}
      \begin{block}{La prioridad}
        La \textbf{calidad de la ingesta}. Si el OCR o las tablas están mal
        extraídos, el mejor modelo de lenguaje responderá mal.
      \end{block}
      \vspace{1mm}
      \begin{block}{Alcance}
        Solo esta norma. Fuera de alcance: diagnóstico, recomendaciones clínicas,
        otras normas, datos de pacientes, interfaz web y \emph{fine-tuning}.
      \end{block}
    \end{column}
  \end{columns}
\end{frame}

% ---------------------------------------------------------------- Arquitectura general
\begin{frame}{Arquitectura general}
  \centering
  \begin{tikzpicture}[>=Stealth]
    % Banda offline
    \fill[grisclaro, rounded corners=6pt] (-0.25,-0.75) rectangle (14.0,0.9);
    \node[anchor=west, font=\scriptsize\bfseries, text=azul] at (-0.1,0.7)
      {ETAPA OFFLINE (una vez)};

    \node[caja] (pdf)  at (1.35,0)   {\texttt{pdf/}\\PDF escaneados};
    \node[caja] (ocr)  at (4.55,0)   {OCR\\+ Docling};
    \node[caja] (proc) at (7.75,0)   {Limpieza\\+ chunks};
    \node[caja] (idx)  at (10.95,0)  {Embeddings\\+ ChromaDB};
    \node[font=\scriptsize, anchor=west, align=left] at (12.6,0) {\texttt{data/}\\\texttt{index/}};

    \draw[flecha] (pdf) -- (ocr);
    \draw[flecha] (ocr) -- (proc);
    \draw[flecha] (proc) -- (idx);

    % Banda online
    \fill[grisclaro, rounded corners=6pt] (-0.25,-3.35) rectangle (14.0,-1.75);
    \node[anchor=west, font=\scriptsize\bfseries, text=naranja] at (-0.1,-1.95)
      {ETAPA ONLINE (por consulta)};

    \node[cajaN] (q)    at (1.35,-2.6)  {Pregunta};
    \node[cajaN] (rec)  at (4.55,-2.6)  {Búsqueda\\híbrida};
    \node[cajaN] (fus)  at (7.75,-2.6)  {Fusión RRF\\top 6--8};
    \node[cajaN] (gen)  at (10.95,-2.6) {LLM\\o extractivo};
    \node[font=\scriptsize, anchor=west, text width=1.3cm, align=left] at (12.6,-2.6)
      {Respuesta\\+ fuentes};

    \draw[flechaN] (q) -- (rec);
    \draw[flechaN] (rec) -- (fus);
    \draw[flechaN] (fus) -- (gen);
    \draw[flechaN] (gen) -- ++(0.75,0);
  \end{tikzpicture}

  \vspace{2mm}
  \footnotesize
  Ingesta local determinista, una sola vez \textbullet{} Consulta en local o en un
  contenedor ligero \textbullet{} Cada fragmento conserva documento, página y sección.
\end{frame}

% ---------------------------------------------------------------- Ingesta
\begin{frame}{Arquitectura: ingesta offline}
  \centering
  \begin{tikzpicture}[>=Stealth]
    \node[arch] (pdf) at (0,0)      {\texttt{pdf/*.pdf}\\2 partes};
    \node[caja] (ocr) at (2.60,0)   {\textbf{1} Docling\\layout + OCR\\+ tablas};
    \node[caja] (doc) at (5.20,0)   {\textbf{2} Fusión\\inter-página\\+ léxico};
    \node[caja] (lim) at (7.80,0)   {\textbf{3} Limpieza\\normalización\\de unidades};
    \node[caja] (chk) at (10.40,0)  {\textbf{4} Chunks\\con metadatos};
    \node[caja] (ix)  at (13.00,0)  {\textbf{5} Índice\\Embeddings +\\ChromaDB};

    \draw[flecha] (pdf) -- (ocr);
    \draw[flecha] (ocr) -- (doc);
    \draw[flecha] (doc) -- (lim);
    \draw[flecha] (lim) -- (chk);
    \draw[flecha] (chk) -- (ix);
  \end{tikzpicture}

  \vspace{4mm}
  \begin{itemize}
    \item \textbf{OCR:} \texttt{tessdata\_best}, español, 300 ppp, motor LSTM; endereza
          páginas giradas y omite formularios en blanco.
    \item \textbf{Extracción:} Docling reconstruye layout, secciones y tablas a
          Markdown/JSON.
    \item \textbf{Trazabilidad:} cada chunk guarda documento, página, sección, vacuna y tipo.
    \item \textbf{Fuente de verdad:} \texttt{data/chunks.jsonl} permite reconstruir el índice
          sin repetir el OCR.
  \end{itemize}
\end{frame}

% ---------------------------------------------------------------- Consulta
\begin{frame}{Arquitectura: consulta online}
  \centering
  \begin{tikzpicture}[>=Stealth]
    \node[cajaN] (q)    at (0,0)      {Pregunta};
    \node[caja]  (emb)  at (2.60,0)   {Embedding\\384 dim};
    \node[cajaN] (vec)  at (5.20,1.0) {ChromaDB\\vectorial top-10};
    \node[cajaN] (bm)   at (5.20,-1.0){BM25\\léxica top-10};
    \node[caja]  (rrf)  at (7.80,0)   {Fusión RRF\\top 6--8};
    \node[caja]  (gen)  at (10.40,0)  {LLM opcional\\o extractivo};
    \node[caja]  (res)  at (13.00,0)  {Respuesta\\+ fuentes};

    \draw[flechaN] (q) -- (emb);
    \draw[flecha] (emb) -- (vec);
    \draw[flecha] (emb) -- (bm);
    \draw[flecha] (vec) -- (rrf);
    \draw[flecha] (bm) -- (rrf);
    \draw[flecha] (rrf) -- (gen);
    \draw[flecha] (gen) -- (res);
  \end{tikzpicture}

  \vspace{3mm}
  \begin{itemize}
    \item \textbf{Recuperación híbrida:} vectores (semántica) + BM25 (términos exactos),
          combinados con \emph{Reciprocal Rank Fusion}.
    \item \textbf{Generación:} LLM con contexto citado mediante API compatible con OpenAI
          (Ollama, OpenAI, etc.).
    \item \textbf{Modo extractivo:} sin LLM configurado, devuelve los fragmentos más
          relevantes con su cita.
    \item \textbf{Abstención:} si la respuesta no está en la norma, lo indica en lugar de
          inventarla.
  \end{itemize}
\end{frame}

% ---------------------------------------------------------------- Stack
\begin{frame}{Stack tecnológico}
  \small
  \begin{columns}[T]
    \begin{column}{0.50\textwidth}
      \begin{itemize}
        \item Python 3.12 y \texttt{uv}
        \item FastAPI + Uvicorn
        \item ChromaDB 1.5 + BM25 en memoria
        \item FastEmbed + ONNX Runtime
              (\texttt{paraphrase-\allowbreak multilingual-\allowbreak MiniLM-L12-v2}, 384 dim.)
      \end{itemize}
    \end{column}
    \begin{column}{0.48\textwidth}
      \begin{itemize}
        \item Docling + RapidOCR (ONNX) en una pasada
        \item TableFormer ACCURATE para tablas
        \item Docker Compose (imagen de consulta ligera, $\approx$850 MB)
        \item LLM opcional vía API compatible con OpenAI
      \end{itemize}
    \end{column}
  \end{columns}

  \vspace{2mm}
  \begin{block}{Endpoints}
    \centering
    \setlength{\tabcolsep}{4pt}
    \begin{tabular}{@{}lll@{}}
      \toprule
      Método & Ruta & Función \\
      \midrule
      GET  & \texttt{/salud}    & Estado e índice \\
      POST & \texttt{/consulta} & Pregunta $\rightarrow$ respuesta + fuentes \\
      POST & \texttt{/buscar}   & Solo recuperación, sin LLM \\
      \bottomrule
    \end{tabular}
  \end{block}
\end{frame}

% ---------------------------------------------------------------- Estado y ejemplos
\begin{frame}{Resultados y ejemplos de consultas}
  \small
  \textbf{Estado:} 140 páginas procesadas \textperiodcentered{} \textbf{790 chunks} indexados
  (17 por revisar) \textperiodcentered{} \textbf{11/11} preguntas de evaluación con la fuente
  esperada en el top 5 \textperiodcentered{} API probada en local y Docker.

  \vspace{2mm}
  \begin{itemize}
    \item \textbf{¿Cuál es el esquema de vacunación de la Tdap en gestantes?}\hfill
          {\scriptsize\color{azul}(p. 34, 6.1.1.17)}\\
          Una dosis de Tdap entre las 20 y 36 semanas de gestación, en cada gestación.
    \item \textbf{¿Cuánto tarda en aparecer la anafilaxia asociada a la DPT?}\hfill
          {\scriptsize\color{azul}(p. 98, Anexo 8)}\\
          De 0 a 1 hora.
    \item \textbf{¿Cuál es la finalidad de la norma?}\hfill
          {\scriptsize\color{azul}(p. 1, FINALIDAD)}\\
          Contribuir a disminuir la morbimortalidad causada por enfermedades inmunoprevenibles.
  \end{itemize}

  \vspace{1mm}
  {\footnotesize\textbf{Abstención:} «¿Cuál es el precio de la vacuna contra el VPH?»\\
   $\rightarrow$ No se encontró esa información en la normativa cargada.}
\end{frame}


\end{document}

%%% Local Variables:
%%% mode: LaTeX
%%% TeX-master: t
%%% End:
